\section{\label{sec:intro}Introduction}

For a closed quantum system, thermal equilibrium is generated by a Hamiltonian:
the Gibbs state is an exponential in the inverse temperature. For an open system
at finite coupling, the subsystem does not in general admit an equilibrium
description in terms of its bare Hamiltonian. The operational equilibrium state
of the subsystem is instead inherited from the global canonical ensemble: one
traces out the bath from the total Gibbs state. \emph{The Hamiltonian of mean force} $H_{\mathrm{MF}}$
(HMF) is the operator that reproduces this reduced equilibrium object as a
Gibbs-like state, and provides the standard starting point for strong-coupling
thermodynamics and mean-force Gibbs-state constructions
\cite{campisiFluctuationTheoremArbitrary2009,talknerColloquiumStatisticalMechanics2020,trushechkinOpenQuantumSystem2022,seifertFirstSecondLaw2016}.

Historically, open-system theory prioritized weak-coupling and Markovian
regimes, where reduced equilibrium can often be approximated by a Gibbs state of a renormalised Hamiltonian. At finite coupling the reduced state inherits explicit
temperature dependence and interaction-induced operator content that is not
captured by a simple renormalisation. Coupling-dependent thermodynamic response
features in quantum Brownian motion and related models highlight this
complexity \cite{hanggiFiniteQuantumDissipation2008,ingoldSpecificHeatAnomalies2009}.
The strong-coupling literature consequently treats the mean-force Gibbs state as
a distinct equilibrium object, with operational ramifications for heat and
energy definitions\cite{espositoNatureHeatStrongly2015,rivasStrongCouplingThermodynamics2020}.

The purpose of this paper is methodological. We present an influence-functional
route to the Hamiltonian of mean force which rewrites the reduced equilibrium density as the averaged state of a stochastic dynamical equation. In the special but physically important
case where the bath couples to a conserved system quantity - or equivalently, the
operator coupling the bath to the bare system commutes with the latter's Hamiltonian - the imaginary-time ordering obstruction disappears and the Gaussian average can be evaluated in
closed form. This yields an explicit operator expression for $H_{\mathrm{MF}}$
with temperature dependence determined directly by the bath correlation kernel
(or spectral density). Outside this commuting sector, the same representation
remains exact but typically does not collapse to a compact local generator.

The HMF literature itself is broad but structured. Canonical definitions and
thermodynamic identities have been developed in strong-coupling thermodynamics and
fluctuation-relation work\cite{campisiFluctuationTheoremArbitrary2009,jarzynskiNonequilibriumWorkTheorem2004,seifertFirstSecondLaw2016,talknerColloquiumStatisticalMechanics2020}., while the ``static'' mean-force Gibbs perspective has been consolidated with the ``dynamical'' return-to-equilibrium perspective in open quantum systems \cite{trushechkinOpenQuantumSystem2022}. Operational questions such as measurability and thermodynamic consistency at strong coupling have also been pursued\cite{strasbergMeasurabilityNonequilibriumThermodynamics2020,rivasStrongCouplingThermodynamics2020}.

Exact or controlled evaluations exist in special cases. For commuting (QND)
interactions, the imaginary-time ordering obstruction disappears and the HMF can
be written explicitly\cite{campisiTalknerHanggi2009Solvable}. Quadratic/Gaussian
models (e.g., damped harmonic oscillators) are solvable because Gaussianity is
preserved, leading to closed operator forms\cite{caldeiraQuantumTunnellingDissipative1983a,grabertQuantumBrownianMotion1988,hiltHamiltonianMeanForce2011}.
Finite-dimensional closures such as spin-boson or single-qubit models provide
additional controlled benchmarks\cite{leggettDynamicsDissipativeTwostate1987}.
Beyond these cases, the mean-force Gibbs state is often accessed through
systematic limits: weak and ultrastrong-coupling expressions and show that, in the ultrastrong limit, the mean-force Gibbs state becomes diagonal in the interaction basis rather than the system Hamiltonian basis\cite{cresserWeakUltrastrongCoupling2021a}. Recent work generalizes the HMF framework to finite baths by introducing a pair of quantum Hamiltonians of mean
force that incorporate bath feedback\cite{duGeneralizedHamiltonianMeanforce2025a},
and structural studies further analyze the operator content of the HMF in
extended settings\cite{burkeStructureHamiltonianMean2024}.

Outside solvable models, most approaches are perturbative or numerical.
Weak-coupling/high-temperature expansions yield controlled but limited
series\cite{cresserWeakUltrastrongCoupling2021a}. Imaginary-time path-integral
methods, stochastic representations, and hierarchical-equations techniques can
compute $\rho_Q(\beta)$ directly but typically do not provide a compact operator
form for $H_{\mathrm{MF}}$ \cite{moixEquilibriumreducedDensityMatrix2012,chenRigorousStochasticMatrix2014,makriExploitingClassicalDecoherence2014,tanimuraReducedHierarchicalEquations2014,songCalculationCorrelatedInitial2015}.
Stochastic Liouville and partition-free approaches provide complementary
numerical access to open-system equilibration without yielding closed-form HMFs
\cite{stockburgerSimulatingSpinbosonDynamics2004,mccaulPartitionfreeApproachOpen2017c}.
These methods are indispensable for quantitative predictions but leave open a
basic representational issue: when does the reduced equilibrium operator admit a
simple generator within a restricted operator family (few-body, spatially local,
or a prescribed algebra)?

The influence-functional formalism is a natural language for this problem.
For Gaussian baths with linear coupling it provides an exact route to
integrating out bath degrees of freedom\cite{feynmanTheoryGeneralQuantum1963a,caldeiraQuantumTunnellingDissipative1983a,grabertQuantumBrownianMotion1988}.
In equilibrium it becomes a Euclidean (imaginary-time) influence functional,
usually bilocal in $\tau$, and admits a Hubbard - Stratonovich rewriting as a
quenched Gaussian-field average\cite{hubbardCalculationPartitionFunctions1959a,stratonovich1957QDistro,stockburgerExactNumberRepresentation2002}.
Related path-integral derivations of quantum Langevin dynamics make explicit the
noise and dissipation structure inherited from the bath\cite{kleinertQuantumLangevinEquation1995,vankampenDerivationQuantumLangevin1997}.
This formulation also clarifies terminology: ``nonlocal'' initially refers to
imaginary-time nonlocality of the kernel~\cite{feynmanTheoryGeneralQuantum1963a,grabertQuantumBrownianMotion1988}, whereas the locality of the HMF is a
question about operator structure on the system Hilbert space~\cite{burkeStructureHamiltonianMean2024}.

The existence of $H_{\mathrm{MF}}$ is not the issue - it is defined by
a logarithm of a traced exponential. The practical obstruction is that the
Euclidean influence functional is nonlocal in imaginary time and, for generic
noncommuting couplings, produces operator content that is difficult to compress
into a small ansatz. The commuting/quantum non-demolition (QND) sector isolates the point at which this obstruction vanishes: time ordering becomes redundant and the Gaussian average
can be carried out in closed form. 

Specialising to commuting couplings is not merely a technical simplification.
Physically, it corresponds to situations in which the environment monitors or
responds to a conserved quantity of the bare system, so that energy exchange is
suppressed and the interaction induces dephasing rather than
transitions~\cite{breuerTheoryOpenQuantum2002}. This
includes quantum nondemolition (QND) measurements, longitudinal ($1/f$) noise in qubits~\cite{makhlinQuantumStateEngineering2001,palmaQuantumComputersDissipation1996},
charge or flux noise coupling to conserved operators, and a wide class of
low-temperature solid-state and mesoscopic settings where relaxation channels
are parametrically weaker than dephasing channels. This "pure dephasing"
sector~\cite{clerkIntroductionQuantumNoise2010} forms a standard chapter in open
quantum systems but is rarely analysed through the specific lens of equilibrium
mean-force operators.

From a formal standpoint, the commuting sector isolates the precise obstruction
to representing the Hamiltonian of mean force as a compact operator. When the
coupling operator commutes with the system Hamiltonian, imaginary-time ordering
becomes redundant and the Euclidean influence functional collapses to a scalar
Gaussian average. The Hamiltonian of mean force can then be obtained exactly by
resumming the Gaussian fluctuations, yielding a closed expression whose
temperature dependence is determined directly by the bath correlation kernel.
This case therefore serves as a sharp benchmark~\cite{campisiTalknerHanggi2009Solvable}: it is the point at which the
influence-functional representation becomes fully evaluable rather than merely
formal.

Methodologically, the commuting sector provides a clean reference against which
more general constructions can be judged. Any approximation scheme or reaction-
coordinate embedding intended to capture strong-coupling equilibrium effects
should reproduce the commuting result exactly when the coupling operator is a
constant of motion. In this sense, the present calculation is not only a solvable
example but a consistency check on broader approaches to mean-force
thermodynamics.


In this paper we give an exact influence-functional and quenched-field
representation of the reduced equilibrium operator for Gaussian baths and derive
an \emph{exact} closed operator expression in the commuting sector. The resulting
Hamiltonian of mean force is fully determined by the bath kernel and can be
checked directly against finite-bath trace-out simulations. This isolates a
clean benchmark regime in which the balance between bath spectral input and
system operator structure is fully explicit.

The paper is structured as follows. In Sec.~\ref{sec:model} we define the total
Hamiltonian and reduced equilibrium objects. Sec.~\ref{sec:quenched} develops
the path-integral influence-functional representation and its
Hubbard - Stratonovich form. Sec.~\ref{sec:influence_to_hmf} derives the exact
commuting Gaussian mean-force Hamiltonian. Sec.~\ref{sec:beyond_gaussian}
briefly records the scope boundary for this paper and defers non-Gaussian and
non-commuting extensions to follow-up work. Sec.~\ref{sec:numerical} presents
paper-grade numerical validation against explicit ED trace-out results, and
Sec.~\ref{sec:discussion} summarizes implications and next steps.

